\chapter{Mô hình đe dọa và ranh giới tin cậy}
\label{chap:threat-model}

\section{Tài sản và mục tiêu bảo vệ}

Các tài sản chính gồm danh tính và phiên đăng nhập; quan hệ user--role--permission;
trạng thái và mã tham chiếu giao dịch; audit trail; log đầu vào; cấu hình rule;
và artifact cảnh báo. Mục tiêu bảo mật là ngăn người dùng thực hiện nghiệp vụ
không được cấp quyền, ngăn một Teller tự hoàn tất giao dịch trong chế độ RBAC,
giữ đủ bằng chứng để giải thích quyết định và không để dữ liệu log không tin
cậy làm hỏng toàn bộ một lần phân tích.

Đối thủ được giả định có thể gửi request HTTP trực tiếp, bỏ qua giao diện,
đoán mật khẩu, chèn chuỗi bất thường vào tham số, dùng tài khoản hợp lệ nhưng
thiếu quyền, hoặc cung cấp file log có dòng sai định dạng. Đối thủ không được
giả định đã kiểm soát máy chủ FastAPI, tiến trình phân tích hay file SQLite.
Nếu các thành phần này bị chiếm quyền, audit log cục bộ và artifact không còn
đảm bảo tính toàn vẹn; PoC chưa có append-only storage, chữ ký số hay remote
attestation để chống lại kịch bản đó.

\section{Ranh giới tin cậy và luồng dữ liệu}

Hình~\ref{fig:trust-boundaries} chỉ ra ba vùng. Trình duyệt và file log nằm
ngoài vùng tin cậy; mọi token, URL, field CSV/JSON và request text đều phải
được xem là dữ liệu không tin cậy. FastAPI và pipeline là vùng thực thi chính
sách. SQLite cùng artifact là vùng lưu trữ cục bộ: chỉ được tin cậy khi quyền
hệ điều hành và máy chủ chưa bị phá vỡ. Ranh giới này giải thích vì sao ẩn nút
ở frontend không phải kiểm soát quyền, và vì sao parser phải xác thực schema
trước khi rule xử lý dữ liệu.

\begin{figure}[htbp]
\centering
\resizebox{\textwidth}{!}{\begin{tikzpicture}[
  node distance=8mm and 12mm,
  box/.style={draw, rounded corners, align=center, minimum height=10mm, minimum width=26mm, fill=blue!6},
  store/.style={draw, cylinder, shape border rotate=90, aspect=0.28, align=center, minimum height=13mm, minimum width=18mm, fill=orange!14},
  boundary/.style={draw, dashed, rounded corners, inner sep=6mm},
  flow/.style={-Latex, thick}
]
\node[box] (browser) {Trình duyệt /\ HTTP client};
\node[box, below=of browser] (log) {CSV/JSON log\ không tin cậy};
\node[box, right=22mm of browser] (api) {FastAPI\ xác thực + RBAC};
\node[box, below=of api] (engine) {Parser + rules\ scoring + correlation};
\node[store, right=22mm of api] (db) {SQLite\ nghiệp vụ + audit};
\node[store, below=of db] (artifact) {Artifact\ CSV/JSON};
\node[boundary, fit=(browser)(log), label={[font=\scriptsize]above:Vùng bên ngoài}] {};
\node[boundary, fit=(api)(engine), label={[font=\scriptsize]above:Vùng thực thi chính sách}] {};
\node[boundary, fit=(db)(artifact), label={[font=\scriptsize]above:Vùng lưu trữ cục bộ}] {};
\draw[flow] (browser) -- node[above, font=\scriptsize]{request/token} (api);
\draw[flow] (api) -- node[above, font=\scriptsize]{transaction/audit} (db);
\draw[flow] (log) -- node[below, font=\scriptsize]{validate} (engine);
\draw[flow] (db) -- node[right, font=\scriptsize]{adapter} (artifact);
\draw[flow] (artifact) -- node[below, font=\scriptsize]{Event} (engine);
\end{tikzpicture}
}
\caption{Ranh giới tin cậy và các điểm kiểm soát của nguyên mẫu.}
\label{fig:trust-boundaries}
\end{figure}

Audit logging cần ghi nhận sự kiện xác thực, thất bại phân quyền và thay đổi
chính sách, đồng thời tránh ghi mật khẩu, token phiên và dữ liệu ngân hàng nhạy
cảm vào log \citep{OWASPLogging}. Bởi vậy adapter Nova Bank chỉ xuất các field
cần cho phân tích: thời điểm, actor, resource, action và outcome. Nó không xuất
password hay token; trường IP dùng giá trị nguồn cục bộ khi audit schema hiện
tại chưa lưu địa chỉ client. SQL Injection tiếp tục được phân tích từ log
gateway/request riêng, vì đưa nguyên payload vào audit nghiệp vụ sẽ làm tăng
rủi ro lưu dữ liệu nhạy cảm và log injection.

\section{Ánh xạ đe dọa sang kiểm soát}

\begin{table}[htbp]
\centering
\small
\begin{tabularx}{\textwidth}{p{0.25\textwidth}p{0.28\textwidth}X}
\toprule
Đe dọa & Kiểm soát hiện có & Phần dư rủi ro \\
\midrule
Bỏ qua giao diện để phê duyệt & Bearer token, permission dependency và policy tại FastAPI & Token demo và SQLite chưa phù hợp triển khai phân tán. \\
Teller tự tạo và tự duyệt & Separation of duties; Controller giữ quyền approve & Admin vẫn có thể đổi policy; cần governance và dual control cho thay đổi chính sách. \\
Đoán mật khẩu & Audit đăng nhập lỗi, luật theo user--IP/cửa sổ & Có thể né ngưỡng bằng password spraying, đổi IP hoặc low-and-slow. \\
SQL Injection & Prepared query ở repository; luật giám sát request log & Regex có thể bỏ sót encoding và không chứng minh truy vấn đã thực thi. \\
Truy cập trái phép & RBAC lookup kết hợp outcome denied & Sai seed quyền hoặc tài khoản máy chủ bị chiếm có thể làm quyết định sai. \\
Log lỗi hoặc độc hại & Schema validation, RowError và output nguyên tử & Chưa có chữ ký, chống sửa đổi, quota hoặc kiểm soát log injection đầy đủ. \\
\bottomrule
\end{tabularx}
\caption{Ánh xạ các đe dọa chính sang kiểm soát và rủi ro còn lại.}
\label{tab:threat-controls}
\end{table}

Mô hình đe dọa giới hạn cách đọc kết quả: detector đánh giá chỉ báo quan sát
được trong log, còn enforcement quyết định một request có được phép hay không.
Một alert SQL Injection không đồng nghĩa dữ liệu đã bị truy xuất; tương tự,
một request bị RBAC từ chối là bằng chứng kiểm soát hoạt động, không phải bằng
chứng tài khoản chắc chắn đã bị xâm nhập.

\section{Giả định an toàn và tiêu chí kiểm chứng}

Mô hình giả định identity trong Event đã được nguồn log gắn chính xác; đồng hồ
các thành phần không lệch đến mức làm sai cửa sổ tương quan; cấu hình và seed
RBAC được quản trị viên có thẩm quyền cung cấp; và đường dẫn file output không
bị tiến trình khác thay thế. Đây là các giả định có thể chấp nhận cho môi trường
đồ án cục bộ nhưng phải trở thành kiểm soát cụ thể khi triển khai: xác thực nguồn
log, đồng bộ thời gian, quản lý thay đổi cấu hình và phân quyền hệ điều hành.

Mỗi ranh giới có một tiêu chí kiểm chứng tương ứng. Ranh giới HTTP được kiểm
tra bằng request trực tiếp tới URL phê duyệt và kỳ vọng HTTP 403, đồng thời giao
dịch vẫn pending. Ranh giới parser được kiểm tra bằng header thiếu, timestamp
sai và loại file không hỗ trợ; dòng lỗi phải được cô lập hoặc lần chạy phải kết
thúc có thông báo kiểm soát. Ranh giới lưu trữ được kiểm tra bằng foreign key,
transaction rollback và phép ghi artifact qua file tạm rồi replace.

Đối với detection, tiêu chí không chỉ là có alert mà còn là khả năng truy ngược
alert tới rule ID, Event ID và evidence. Đối với audit integration, tiêu chí là
chuỗi thao tác thật trong Nova Bank được xuất thành Event hợp lệ và tạo alert
mong đợi mà không cần gắn \texttt{expected\_label}. Các tiêu chí này biến mô
hình đe dọa thành điều kiện kiểm thử; chúng không loại bỏ rủi ro dư đã nêu trong
Bảng~\ref{tab:threat-controls}.
